\section{Introduction}\label{sec:introduction}
\begin{CJK*}{UTF8}{gbsn}

\nextrevision{零样本音频分类旨在利用类别语义识别缺少目标任务训练样本的声音类别，从而减少对音频标注数据的依赖。早期方法通过匹配音频表示与类别语义嵌入，实现对未见类别的识别~\cite{xie2021zero}。对比语言--音频预训练模型（Contrastive Language--Audio Pretraining，CLAP）进一步通过大规模音频--文本对比学习建立共享表示空间，使自然语言描述能够直接用于声音分类~\cite{elizalde2023clap}。后续研究通过扩充训练数据、改进预训练目标及音频建模方式，持续增强音频--语言表示与零样本泛化能力~\cite{wu2023clap,niizumi2024m2dclap,mei2026slap}。然而，当类别主要由简短名称表示时，声音的来源、发生场景及其与其他事件的语义联系难以得到充分表达，限制了分类过程中可利用的判别信息。}

\nextrevision{为丰富类别的语义表示，已有研究从文本提示和结构化知识两个方向展开探索。ReCLAP通过改写音频描述和生成类别增强提示，补充声音特征与场景信息；PAT根据输入音频选择和组合提示，并调整音频表示，以改善音频与文本的匹配~\cite{ghosh2024reclap,seth2025pat}。这些工作从不同角度说明了文本内容及其与输入音频的适配对零样本分类的重要性。结构化知识则提供了另一种补充语义信息的途径。iKnow-audio构建音频中心知识图谱（Audio-centric Knowledge Graph，AKG），通过经过筛选的关系集合检索与候选类别相关的尾实体，并将其与类别名称组合为增强提示，用于修正CLAP的候选排序~\cite{olvera-etal-2025-iknow}。该研究展示了将声音事件之间的结构化联系引入零样本音频分类的可行性，也进一步引出了如何使图谱知识适配具体音频的问题。}

\nextrevision{尽管上述研究丰富了零样本音频分类可利用的语义信息，图谱知识在具体样本上的选择、表达与融合仍有进一步改进的空间。首先，经过筛选的关系集合能够提供类别相关知识，但同一关系对不同音频样本的判别价值可能不同，仍需根据当前音频确定参与知识增强的关系。其次，iKnow-audio通过拼接类别名称与检索到的尾实体构造增强提示，未显式保留连接二者的关系语义，因而需要一种能够表达实体间关系、同时约束额外内容扩写的文本构造方式。最后，在分数融合方面，原始类别分数与增强提示分数被共同聚合，知识证据的汇总与其对基础预测的影响在同一聚合过程中确定，缺少对这两个环节的分别控制。当某一类别的知识提示数量较多或匹配分数较高时，知识证据可能对聚合结果产生过大影响，使原始CLAP预测的作用难以单独调节。因此，如何依据当前音频选择关系、完整表达所选知识的关系语义，并分别组织知识证据聚合与原始预测融合，是本文关注的问题。}

\nextrevision{针对上述问题，本文提出样本自适应知识整合框架（Sample-Adaptive Knowledge Integration，SAKI），在保持CLAP参数不变的条件下，利用结构化知识改善候选类别排序。在离线阶段，音频对齐知识语言化模块（Audio-Aligned Knowledge Verbalization，AAKV）引导大语言模型在保留实体及关系语义的前提下，将检索得到的三元组改写为以声音为主体的自然语言描述，并缓存对应的文本向量。在在线阶段，SAKI为每种关系构建一个利用该关系知识进行分类的关系专家，依据专家间的预测共识与各专家的分类间隔，选择参与当前音频分类的关系。随后，框架聚合所选关系提供的知识证据，并与原始CLAP分数进行分层融合，完成候选类别重排序。该框架无需利用目标数据集的训练样本训练关系选择器，在线推理过程不使用测试标签。}

本文的主要贡献如下：
\begin{itemize}
  \item \nextrevision{提出样本级关系选择方法，利用关系专家的预测信息，使参与知识增强的关系随输入音频调整。}
  \item \nextrevision{提出AAKV，在保留实体与关系语义的基础上，将图谱三元组转化为面向音频匹配的文本证据。}
  \item \nextrevision{构建分层知识融合机制，协调知识证据与原始CLAP预测，并通过跨数据集对照与消融实验评估完整框架及各模块的作用。}
\end{itemize}

\nextrevision{本文在五个公开音频数据集上开展评估。相较于受控复现的iKnow-audio，SAKI的Hit@1提高了2.08--12.73个百分点，MRR提高了1.02--6.68个百分点。组合消融结果显示，完整配置取得最高的五数据集等权平均Hit@1和MRR；关系选择与知识语言化对照进一步检验了相应模块的作用。此外，本文通过统计分析、逐样本预测变化、参数敏感性及计算成本分析，考察方法的性能收益与适用边界。}

\nextrevision{本文其余部分安排如下：第2节介绍相关研究，第3节阐述SAKI的总体框架及三个核心模块，第4节报告实验设置、主要结果与进一步分析，第5节讨论主要发现与研究局限，第6节总结全文。}

\end{CJK*}
